\documentclass[11pt,a4paper]{article}

\usepackage[utf8]{inputenc}
\usepackage[T1]{fontenc}
\usepackage{geometry}
\geometry{margin=2.2cm}
\usepackage{hyperref}
\usepackage{booktabs}
\usepackage{graphicx}
\usepackage{listings}
\usepackage{xcolor}
\usepackage{amsmath}
\usepackage{enumitem}

\hypersetup{
    colorlinks=true,
    linkcolor=blue!70!black,
    urlcolor=blue!70!black,
    citecolor=blue!70!black
}

\lstset{
    basicstyle=\ttfamily\small,
    backgroundcolor=\color{gray!8},
    frame=single,
    framerule=0.4pt,
    rulecolor=\color{gray!40},
    breaklines=true,
    columns=fullflexible,
    keepspaces=true,
    showstringspaces=false,
    xleftmargin=4pt,
    xrightmargin=4pt,
    aboveskip=6pt,
    belowskip=6pt
}

\title{\textbf{Sentence Splitter}\\[4pt]
\large Fine-tuning XLM-RoBERTa-large for Sentence Boundary Detection}
\author{Luca Tam\\
\small Sapienza University of Rome\\
\small Multilingual Natural Language Processing}
\date{}

\begin{document}
\maketitle

\section{Introduction}

This report describes a sentence boundary detection system built by fine-tuning
\textbf{XLM-RoBERTa-large} (560M parameters) as a per-token binary classifier.
The model is trained on 8 Universal Dependencies treebanks (4 English, 4 Italian)
and achieves a macro F1 of \textbf{0.9863}, improving over the NLTK Punkt
baseline (0.9411) by +4.5\% and spaCy (0.9519) by +3.4\% absolute.

The full source code, training scripts, and evaluation pipeline are available at:

\begin{center}
\url{https://github.com/LucaTamSapienza/sentence_splitter}
\end{center}

\section{Project Structure}

\begin{lstlisting}
sentence_splitter/
  src/
    model.py          # Model architecture
    data.py           # Data parsing, windowing, dataset classes
    train.py          # Training loop (fp16, early stopping)
    inference.py      # Inference pipeline (sliding windows)
    evaluate.py       # Evaluation metrics (F1, exact match)
    ensemble.py       # Multi-model weighted blending
    rules.py          # Rule-based boundary signals
    utils.py          # Config loading, seeding, logging
    train_xlmr.py     # CLI: training script
    evaluate_xlmr.py  # CLI: threshold tuning + evaluation
    predict.py        # CLI: inference on arbitrary text
    run_baselines.py  # CLI: NLTK + spaCy baseline comparison
    optimize.py       # CLI: ensemble weight grid search
  configs/
    xlmr.yaml         # Hyperparameter configuration
  UD_*/               # 8 Universal Dependencies treebanks
  checkpoints/        # Saved model weights
  requirements.txt
\end{lstlisting}

\section{Method}

\subsection{Architecture}

The model applies a linear classification head on top of each token representation
produced by XLM-RoBERTa-large:
%
\[
  \hat{y}_i = \sigma\!\bigl(\mathbf{w}^\top \operatorname{Dropout}(\mathbf{h}_i) + b\bigr),
  \qquad \mathbf{h}_i \in \mathbb{R}^{1024}
\]
%
where $\mathbf{h}_i$ is the hidden state of the $i$-th subtoken, $\sigma$ is the
sigmoid function, and $\hat{y}_i$ is the probability that token $i$ ends a sentence.

\subsection{Sliding Windows}

Documents are split into overlapping windows of 510 content tokens (plus
\texttt{[CLS]}/\texttt{[SEP]}) with a stride of 256 tokens (50\% overlap).
Predictions in overlapping regions are averaged, providing ensemble-like
robustness without multiple models.

\subsection{Training Details}

\begin{itemize}[nosep]
  \item \textbf{Loss}: Focal binary cross-entropy ($\alpha{=}0.75$, $\gamma{=}2.0$) to
        handle the ${\sim}1{:}25$ boundary-to-non-boundary class imbalance.
  \item \textbf{Optimizer}: AdamW, learning rate $2{\times}10^{-5}$, weight decay $0.01$,
        linear warmup (10\%) followed by linear decay.
  \item \textbf{Batch size}: 16 with gradient accumulation of 2 (effective 32).
  \item \textbf{Precision}: FP16 mixed precision.
  \item \textbf{Early stopping}: patience of 3 epochs on dev boundary F1.
  \item \textbf{Sampling}: Temperature-weighted ($T{=}2.0$) across all 7 training
        treebanks to prevent large corpora from dominating.
  \item \textbf{Seed}: 42 for full reproducibility.
\end{itemize}

\section{Reproducibility Guide}

\subsection{Environment Setup}

\begin{lstlisting}[language=bash]
git clone https://github.com/LucaTamSapienza/sentence_splitter.git
cd sentence_splitter
pip install -r requirements.txt
\end{lstlisting}

For GPU support with CUDA 12.1:

\begin{lstlisting}[language=bash]
pip install torch --index-url https://download.pytorch.org/whl/cu121
\end{lstlisting}

To run the NLTK and spaCy baselines, also execute:

\begin{lstlisting}[language=bash]
python -m nltk.downloader punkt_tab
python -m spacy download en_core_web_sm
python -m spacy download it_core_news_sm
\end{lstlisting}

\subsection{Training}

The training data consists of \texttt{.sent\_split} files from 8 Universal
Dependencies treebanks (included in the repository under \texttt{UD\_*/}), where
sentence boundaries are marked with \texttt{<EOS>} tokens.

\begin{lstlisting}[language=bash]
python src/train_xlmr.py \
    --data_dir ./ \
    --output_dir checkpoints \
    --model_name_or_path xlm-roberta-large
\end{lstlisting}

The best checkpoint is saved to \texttt{checkpoints/best\_xlmr\_model.pt}.
Training takes approximately 1--2 hours on a single GPU (e.g., NVIDIA A100).

\subsection{Evaluation and Threshold Tuning}

After training, evaluate on the dev/test splits and sweep thresholds:

\begin{lstlisting}[language=bash]
python src/evaluate_xlmr.py \
    --data_dir ./ \
    --model_path checkpoints/best_xlmr_model.pt \
    --model_name_or_path xlm-roberta-large \
    --predictions_dir outputs/predictions
\end{lstlisting}

The script sweeps thresholds on dev and selects the best one (typically $t{=}0.65$),
then evaluates on all test sets.

\subsection{Inference on New Text}

To split sentences in an arbitrary text file:

\begin{lstlisting}[language=bash]
python src/predict.py \
    --input_file path/to/input.txt \
    --model_path checkpoints/best_xlmr_model.pt \
    --model_name_or_path xlm-roberta-large \
    --threshold 0.65 \
    --output_file outputs/predicted.txt
\end{lstlisting}

The output lists numbered sentences:

\begin{lstlisting}
1) This is the first sentence.

2) This is the second sentence.
\end{lstlisting}

\subsection{Running Baselines}

To reproduce the NLTK and spaCy baselines for comparison:

\begin{lstlisting}[language=bash]
python src/run_baselines.py --split test --verbose
\end{lstlisting}

\section{Results}

Table~\ref{tab:results} reports boundary-level F1 scores on all 8 test sets.

\begin{table}[h]
\centering
\small
\begin{tabular}{lccc}
\toprule
\textbf{Dataset} & \textbf{NLTK} & \textbf{spaCy} & \textbf{XLM-R} ($t{=}0.65$) \\
\midrule
EN-EWT     & 0.803 & 0.819 & 0.957 \\
EN-GUM     & 0.919 & 0.923 & 0.979 \\
EN-ParTUT  & 0.993 & 0.987 & 0.997 \\
EN-PUD     & 0.991 & 0.997 & 0.988 \\
IT-ISDT    & 0.954 & 0.991 & 0.997 \\
IT-MarkIT  & 0.917 & 0.984 & 0.997 \\
IT-ParTUT  & 0.993 & 0.993 & 1.000 \\
IT-VIT     & 0.958 & 0.921 & 0.976 \\
\midrule
\textbf{Macro F1} & \textbf{0.9411} & \textbf{0.9519} & \textbf{0.9863} \\
\bottomrule
\end{tabular}
\caption{Boundary-level F1 on all 8 UD test sets. Threshold tuned on dev ($t{=}0.65$).}
\label{tab:results}
\end{table}

The fine-tuned XLM-RoBERTa model outperforms both baselines by a
significant margin across nearly all datasets, with the largest gains on harder corpora
such as EN-EWT (+15.4\% over NLTK) and IT-VIT (+5.5\% over spaCy). The model
handles both English and Italian with a single multilingual checkpoint.

\subsection{Real-World Gold Test Sets}

Beyond the UD benchmarks, we manually created gold-annotated test sets to evaluate
the model on realistic, challenging text that standard treebanks do not cover.
These are available in the \texttt{data/} directory as JSONL files built with
\texttt{src/build\_test.py}. They can be evaluated with:

\begin{lstlisting}[language=bash]
python src/eval_test.py \
    --input data/full_test.jsonl \
    --checkpoint checkpoints/best_xlmr_model.pt \
    --threshold 0.65
python src/eval_test.py \
    --input data/Fra.jsonl \
    --checkpoint checkpoints/best_xlmr_model.pt \
    --threshold 0.65
\end{lstlisting}

The test sets include:

\begin{itemize}[nosep]
  \item \textbf{hard\_test} (20 sentences): Edge cases such as abbreviations
        (Dr., Ph.D., U.S., S.p.A., art.), ellipses, multi-line sentences,
        numbered lists, and headers without punctuation.
  \item \textbf{extended\_test} (62 sentences): Academic and technical language,
        URLs, mathematical notation (O($n \log n$), $f(x)$), Italian legal
        references (d.lgs., Corte di Cassazione), citations (Smith et al., 2023),
        and informal conversational text.
  \item \textbf{Fra} (165 sentences): A comprehensive real-world test set covering
        Italian legal prose (sentenze, ordinanze, T.A.R., Codice Privacy),
        medical/scientific text (E.\ coli, temperature conversions), financial
        language (fatturato, EUR/USD), system logs, regex patterns, bibliographic
        citations, nested parenthetical structures, and mixed Italian/English content.
\end{itemize}

\begin{table}[h]
\centering
\small
\begin{tabular}{lccccc}
\toprule
\textbf{Test set} & \textbf{Sents} & \textbf{P} & \textbf{R} & \textbf{F1} & \textbf{ExMatch} \\
\midrule
hard\_test       &  20  & 1.000 & 0.950 & 0.974 & 0.950 \\
extended\_test   &  62  & 0.938 & 0.984 & 0.961 & 0.894 \\
Fra              & 165  & 0.896 & 0.988 & 0.939 & 0.803 \\
\bottomrule
\end{tabular}
\caption{Results on manually annotated gold test sets ($t{=}0.65$).
ExMatch = fraction of sentences exactly matching the gold segmentation.}
\label{tab:gold}
\end{table}

The model maintains strong performance on these out-of-domain texts.
Most errors involve over-splitting at ambiguous punctuation in complex legal
abbreviations or enumerated lists---cases where even human annotators may
disagree. The high recall (${\geq}0.95$) shows that the model rarely misses
a true sentence boundary.

\section{Conclusion}

Fine-tuning XLM-RoBERTa-large with focal loss, sliding-window inference, and
temperature-weighted multilingual sampling yields a sentence splitter that
substantially outperforms NLTK Punkt and spaCy, achieving 0.9863 macro F1 on
8 Universal Dependencies test sets. The full pipeline: training, evaluation, and
inference is reproducible from the public repository.

\end{document}
